\section{Related Work}
\label{sec:related}

\paragraph{LLM-as-a-Judge.}
Using LLMs to evaluate the outputs of other models has gained widespread adoption as a cost-effective alternative to human annotation. Strong LLMs can approximate human preferences on subjective dimensions~\citep{liu-etal-2023-g}. As a result, they have been adopted for reward modeling and reinforcement learning~\citep{son2024llmasajudgerewardmodel}. More recently, LLM judges have been extended to objective tasks such as factual evaluation. SAFE~\citep{NEURIPS2024_937ae0e8} uses an LLM agent with web search to verify individual facts in long-form responses. Similarly, SAGE~\citep{tale_sb} employs a tool-augmented agent that iteratively retrieves and synthesizes external evidence for reference-free QA evaluation. These works improve judge accuracy but provide no formal guarantee on the reliability of verdicts.

\paragraph{Uncertainty quantification for LLM evaluation.}
Existing methods generally fall into two categories. White-box methods estimate uncertainty from the judge's token probabilities. The most common choice is the predictive entropy of the output distribution~\citep{malinin2021uncertainty}. These probabilities are often well calibrated with respect to correctness~\citep{kadavath2022languagemodelsmostlyknow}. Black-box methods, in contrast, estimate uncertainty by querying the judge repeatedly and measuring how consistent its verdicts are~\citep{002WSLCNCZ23, wang-etal-2024-conu}. Both signals correlate with judge accuracy and can flag low-confidence verdicts~\citep{scope}. However, they remain heuristic. For instance, a threshold chosen by hand or tuned on a validation set gives no formal guarantee on the error rate among accepted verdicts.

\paragraph{Uncertainty-based routing and tool use.}
A parallel line of work uses uncertainty to determine when additional resources are needed. FrugalGPT~\citep{chen2023frugalgptuselargelanguage} and AutoMix~\citep{NEURIPS2024_ecda225c} use confidence to route queries across models with different capabilities, trading off cost against quality. In contrast, FLARE~\citep{jiang-etal-2023-active} uses token-level uncertainty to decide when to retrieve external information during generation. More recently, SAGE~\citep{tale_sb} introduced search-augmented retrieval for LLM judges, where the judge acts as an agent that gathers targeted evidence when evaluating free-form answers. However, these approaches often rely on unconditional retrieval, which introduces unnecessary computation. Furthermore, they do not provide formal guarantees on the error rate of the resulting evaluations.

\paragraph{Risk-controlled selective evaluation.}
Conformal prediction provides distribution-free, finite-sample coverage guarantees by constructing prediction sets calibrated on held-out data~\citep{angelopoulos2023gentle}. Conformal risk control generalizes this beyond coverage to bound arbitrary loss functions, including the FDR, at user-specified levels~\citep{angelopoulos2024conformal}. Recent work has applied these guarantees to question answering. SConU builds prediction sets that cover admissible answers with high probability~\citep{wang-etal-2025-sconu}. Similarly, COIN~\citep{wang2025coin} calibrates an uncertainty threshold that bounds the FDR among accepted answers using Clopper--Pearson upper confidence bounds~\citep{be7c0fd0_cp}. Closer to our setting, Trust or Escalate~\citep{jung2025trust} and SCOPE~\citep{scope} adapt risk-controlled selective prediction to pairwise evaluation and provide guarantee on human agreement among accepted judgments. However, both target subjective preference tasks where the judge compares two responses rather than evaluating factual correctness (i.e., objective evaluation).

Our framework introduces risk-controlled retrieval-augmented evaluation, where retrieval is selectively activated based on calibrated uncertainty thresholds. Unlike unconditional retrieval approaches, it provides formal FDR guarantees for the accepted verdicts.